As stated in the previous chapter, the \gls{SHM} is declared by the fuzzer and the forkserver can 
know the ID of the \gls{SHM} thanks to an environment variable set by the fuzzer.
The forkserver is injected as a constructor function of a \gls{SO} for which a DT\_NEEDED dependency 
has been added to the target binary.
But, as stated before, the dynamic linker executes the constructor of the \gls{SO} (and therefore the forkserver) before 
the actual initialization of the binary.
Moreover, the \gls{SHM} address is saved in a variable of the \gls{SO} and is not, by default, visible inside the target binary.
The best approach that allows a variable to be accessible outside a binary is to export the symbol 
that represents the variable.

As seen in \cref{sec:loading_relocations}, the relocations for external symbols are resolved during the 
execution of the dynamic linker before the main (or at the first invocation of the function in the case of lazy binding).
Since the \gls{SO} contains the symbol holding the address of the \gls{SHM}, the idea to let the target binary 
know the address of the \gls{SHM} is:
\begin{itemize}
    \item The \gls{SO} exports that symbol, declaring it with\\\texttt{\_\_attribute\_\_((visibility("default")))} \footnote{
        It is not mandatory to explicitly write the visibility attribute, because each symbol has ``default'' visibility by default.
    }
    \item Add a relocation entry to the target binary. 
\end{itemize}
By doing this, the dynamic linker, at binary loading, will resolve the relocation, finding the address of the variable 
that stores the address of the \gls{SHM} inside the \gls{SO}, and write it into the \textit{.got} of the target binary.

Since, as stated in \cref{sec:forkserver_insertion}, it is difficult to add a \gls{GOT} entry in the general case (because 
it depends on the presence of available space in the \gls{GOT}), the chosen approach has been to add a new section 
to the binary where the needed symbols, after relocation, will be written.
Now the complete method to make the \gls{SHM} accessible to the target instrumented binary can be reviewed: 
\begin{enumerate}
    \item Add a new section to the binary. This section will be needed as a \gls{GOT} in which the resolved symbol 
        values will be written.
    \item Add a dynamic symbol to the binary. 
        This is the local symbol that will be referenced by the relocation table.
        The name of the local symbol must be the same as the symbol exported by the \gls{SO} that has already been 
        injected as a dependency of the binary.
    \item Add a dynamic relocation entry. This relocation entry must link the symbol to be relocated with the memory address 
        at which the resolved value will be written (i.e. the section previously created to host the \gls{GOT}).
\end{enumerate}
\begin{figure}[ht]
\centering
        \begin{subfigure}[h]{0.8\textwidth}
        \centering
        \includegraphics[width=\linewidth]{SHMDiscoveryA.png}
        \caption{Before relocation}
        \label{fig:shm_relocation1}
    \end{subfigure}
    \vspace{0.5cm}
    \begin{subfigure}[h]{0.8\textwidth}
        \centering
        \includegraphics[width=\linewidth]{SHMDiscoveryB.png}
        \caption{After relocation but before the forkserver execution}
        \label{fig:shm_relocation2}
    \end{subfigure}
    \vspace{0.5cm}
    \begin{subfigure}[h]{0.8\textwidth}
        \centering
        \includegraphics[width=\linewidth]{SHMDiscoveryC.png}
        \caption{After forkserver execution}
        \label{fig:shm_relocation3}
    \end{subfigure}
    \caption{Shared Memory relocation between the target binary and the Shared Object}
    \label{fig:shm_relocation}
\end{figure}
According to \cref{fig:dynamicallyLinkedProgramLoad}, the relocations are resolved before the forkserver invocation (because the 
forkserver lies inside the constructor of the shared object).
The correct value for the \gls{SHM} address is obtained by the forkserver, but the symbol is resolved beforehand, and therefore 
the address of the symbol is written inside the \gls{GOT} before the variable contains the address of the \gls{SHM}.
When the main starts, the forkserver has already been executed, and dereferencing the relocated symbol will result 
in the correct value initialized by the forkserver.
The process can be seen in \cref{fig:shm_relocation}.
In \cref{fig:shm_relocation1} the \gls{GOT} contains the address 0x0, because the relocation has not yet been performed 
(assume this is after the \gls{SO} loading but before symbol relocation).
Just after the relocation of symbols is performed, the \gls{GOT} will contain the address of the external variable, 
but that variable will still contain 0x0, because the \gls{SHM} has not been attached yet, as shown in \cref{fig:shm_relocation2}.
After forkserver execution, the external variable will also contain the address of the \gls{SHM}, as shown in \cref{fig:shm_relocation3}.